\documentclass[11pt]{article}

\usepackage[margin=1in]{geometry}
\usepackage{amsmath,amssymb,amsthm}
\usepackage{tikz-cd}

\newtheorem{theorem}{Theorem}[section]
\newtheorem{proposition}[theorem]{Proposition}
\newtheorem{lemma}[theorem]{Lemma}
\newtheorem{exercise}[theorem]{Exercise}
\theoremstyle{definition}
\newtheorem{definition}[theorem]{Definition}
\newtheorem{example}[theorem]{Example}
\newtheorem{calculation}[theorem]{Standard computation}
\theoremstyle{remark}
\newtheorem{remark}[theorem]{Remark}

\newcommand{\A}{\mathbb{A}}
\newcommand{\C}{\mathbb{C}}
\newcommand{\F}{\mathbb{F}}
\newcommand{\Pj}{\mathbb{P}}
\newcommand{\Q}{\mathbb{Q}}
\newcommand{\Z}{\mathbb{Z}}
\newcommand{\K}{K}
\newcommand{\V}{\operatorname{V}}
\newcommand{\image}{\operatorname{im}}
\newcommand{\Frac}{\operatorname{Frac}}
\newcommand{\Gal}{\operatorname{Gal}}
\newcommand{\sourceunclear}[2]{%
  \par\smallskip
  \noindent\textbf{TODO: unclear from source, page #1.} #2
  \par\smallskip
}

\title{Algebraic Geometry and Galois Theory\\
       \large Lecture Notes from June 6, 2026}
\author{Based on \texttt{Note Jun 6, 2026 algebraic geometry 1.pdf}}
\date{June 6, 2026}

\begin{document}

\maketitle
\tableofcontents

\section{Guiding Theme}

% Source: PDF pp. 1--9
These notes organize the source around one main idea: algebraic geometry turns
geometric objects into fields of functions, and those fields can then be
studied by field-theoretic methods.  In the examples below, projective curves
such as rational curves and elliptic curves give finite extensions of simpler
function fields.  This is the point at which algebraic geometry meets Galois
theory.

The route is:
\[
\begin{tikzcd}[column sep=large]
  \text{variety}
    \arrow[r, "\text{regular functions}"]
  & \text{coordinate ring}
    \arrow[r, "\text{fractions}"]
  & \text{function field}
    \arrow[r, "\text{comparison}"]
  & \text{field extension}.
\end{tikzcd}
\]
The final example is an elliptic curve \(E\) whose function field is a
quadratic extension of \(\C(x)\).

The source is handwritten.  The main mathematical content has been written in
a textbook style, while unreadable arrows, labels, or diagrams are explicitly
marked as unclear.  Several computations below are included as standard
background calculations so that the algebra can be checked line by line; these
are labeled as standard computations rather than as newly recovered source
content.

\section{Learning Goals and Notation}

\subsection*{Prerequisites}

These notes assume basic field theory, vector spaces, polynomial rings, and
quotient rings.  Some standard background from algebraic geometry and Riemann
surface theory is used as context, especially projective space, homogeneous
polynomials, and the fact that compact Riemann surfaces can be realized as
projective algebraic curves.

\subsection*{Learning goals}

After reading these notes, a student should be able to:
\begin{itemize}
\item explain how a finite field extension is a vector space over the base
field;
\item carry out explicit basis, coordinate-ring, fraction-field, and
invariant-field calculations;
\item define affine and projective varieties, coordinate rings, and function
fields;
\item use homogeneous coordinates on projective space;
\item compute the images of the quadratic and cubic Veronese maps;
\item pass from a projective Weierstrass equation to its affine chart;
\item describe why the elliptic curve \(y^2=x^3+ax+b\) gives a quadratic
extension of \(\C(x)\) with Galois group \(\Z/2\Z\).
\end{itemize}

\subsection*{Notation}

\[
\begin{array}{c|l}
\text{Notation} & \text{Meaning}\\ \hline
k & \text{a base field}\\
\C,\Q,\F_q & \text{the complex numbers, rationals, and a finite field}\\
\A^n_k & \text{affine }n\text{-space over }k\\
\Pj^N_k & \text{projective }N\text{-space over }k\\
\V(I) & \text{zero locus of an ideal }I\\
k[X] & \text{coordinate ring of an affine variety }X\\
\K(X) & \text{function field of an irreducible variety }X\\
{[x_0:\cdots:x_N]} & \text{homogeneous coordinates}\\
E & \text{the elliptic curve in Weierstrass form}\\
F\subseteq L & \text{an extension of fields}\\
\Gal(L/F) & \text{Galois group of }L/F
\end{array}
\]

\section{Field Extensions as Linear Algebra}

We begin with the field-theoretic side of the story.  A finite extension of
fields is not just a larger field: it is also a finite-dimensional vector space
over the smaller field.  This is the linear-algebraic structure that later
appears when a map of curves produces an extension of function fields.

% Source: PDF p. 1
\begin{definition}
Let \(F \subseteq E\) be fields.  The extension \(E/F\) is called
\emph{finite} if \(E\) is finite-dimensional as a vector space over \(F\):
\[
  [E:F]:=\dim_F E<\infty.
\]
It is \emph{finitely generated} if there exist
\(\alpha_1,\ldots,\alpha_n\in E\) such that
\[
  E=F(\alpha_1,\ldots,\alpha_n).
\]
\end{definition}

% Source: PDF p. 1
\begin{remark}
The source emphasizes that finite field extensions can be studied ``via
linear algebra.''  If \(F\subseteq E\), then \(E\) is an \(F\)-vector space:
addition is the addition of the field \(E\), and scalar multiplication is
\[
  F\times E\longrightarrow E,\qquad (a,v)\longmapsto av.
\]
The multiplication on \(E\) is compatible with this vector-space structure:
it is \(F\)-bilinear as a map
\[
  E\times E\longrightarrow E,\qquad (u,v)\longmapsto uv.
\]
Thus a finite extension is simultaneously a field and a finite-dimensional
vector space over the smaller field.
\end{remark}

\begin{calculation}[Two extensions as vector spaces]
The following computations are standard background examples of the slogan
``field extensions can be studied by linear algebra.''

First consider \(\C/\mathbb{R}\).  As a real vector space,
\[
  \C=\mathbb{R}\cdot 1\oplus \mathbb{R}\cdot i,
  \qquad [\C:\mathbb{R}]=2.
\]
If \(z=a+bi\), then its coordinate column in the ordered basis \((1,i)\) is
\[
  [z]_{(1,i)}=\begin{pmatrix}a\\ b\end{pmatrix}.
\]
Multiplication by a fixed element \(u=p+qi\) is an \(\mathbb{R}\)-linear map
\[
  m_u\colon \C\longrightarrow \C,\qquad z\longmapsto uz.
\]
Since
\[
  (p+qi)(a+bi)=(pa-qb)+(qa+pb)i,
\]
the matrix of \(m_u\) in the basis \((1,i)\) is
\[
  [m_u]_{(1,i)}
  =
  \begin{pmatrix}
    p & -q\\
    q & p
  \end{pmatrix}.
\]
Thus field multiplication becomes an explicit linear transformation.

Similarly, let \(d\in\Q\) be squarefree and not a square.  Then
\[
  \Q(\sqrt d)=\Q\cdot 1\oplus \Q\cdot \sqrt d,
  \qquad [\Q(\sqrt d):\Q]=2.
\]
For \(u=a+b\sqrt d\), multiplication by \(u\) is determined by
\[
  (a+b\sqrt d)(r+s\sqrt d)=(ar+bds)+(as+br)\sqrt d.
\]
Hence, in the ordered basis \((1,\sqrt d)\),
\[
  [m_u]_{(1,\sqrt d)}
  =
  \begin{pmatrix}
    a & bd\\
    b & a
  \end{pmatrix}.
\]
\end{calculation}

\begin{remark}[Section summary]
A field extension \(E/F\) can be studied by linear algebra when \(E\) is
finite-dimensional over \(F\).  Later, a geometric map of curves will produce
such an extension at the level of function fields.
\end{remark}

\sourceunclear{1}{The handwritten operation diagram compares
\((F,+_F,\cdot_F)\) and \((E,+_E,\cdot_E)\), but several arrows and labels are
not legible enough to transcribe with confidence.}

\section{Varieties and Function Fields}

The geometric objects in these notes are varieties: spaces cut out by
polynomial equations.  The algebra attached to such a space is its coordinate
ring, and the field of fractions of that ring is its function field.

% Source: PDF p. 2
\begin{definition}
Let \(k\) be a field.  An affine algebraic set in \(\A^n_k\) is the common
zero locus
\[
  \V(I)=\{p\in \A^n_k \mid f(p)=0\text{ for every }f\in I\}
\]
of an ideal \(I\subseteq k[x_1,\ldots,x_n]\).  In these notes, an affine
variety is an algebraic set whose vanishing ideal is prime.  Equivalently, its
coordinate ring is an integral domain.
\end{definition}

% Source: PDF p. 2
\begin{definition}
A smooth \(n\)-manifold is a topological space \(M\) which is locally modeled
on Euclidean space.  More precisely, for every \(p\in M\) there is an open
neighborhood \(U_p\subseteq M\) and a homeomorphism
\[
  U_p\cong \mathbb{R}^n
\]
compatible with a differentiable atlas.
\end{definition}

% Source: PDF p. 2
\begin{remark}
The contrast in the source is the following.  A smooth manifold is locally
modeled on \(\mathbb{R}^n\), whereas an algebraic variety is locally or
globally described by polynomial equations.  The geometric intuition can look
similar, but the definitions impose different kinds of structure.
\end{remark}

% Source: PDF p. 2
\begin{definition}
Let \(X=\V(I)\subseteq \A^n_k\) be an affine variety with \(I\) prime.  Its
coordinate ring is
\[
  k[X]=k[x_1,\ldots,x_n]/I.
\]
Since \(I\) is prime, \(k[X]\) is an integral domain.  Its \emph{function
field} is
\[
  \K(X):=\Frac(k[X]).
\]
\end{definition}

\begin{remark}
Elements of \(\K(X)\) are rational functions on \(X\).  They can be thought of
as quotients \(f/g\) of polynomial functions, with \(g\) not identically zero.
This is the algebraic-geometric analogue of passing from a space to the field
of functions on it.
\end{remark}

\begin{example}
For the affine line \(X=\A^1_{\C}\), the coordinate ring is \(\C[x]\), so
\[
  \K(\A^1_{\C})=\Frac(\C[x])=\C(x).
\]
Thus \(\C(x)\) is the field of rational functions in one variable.
\end{example}

\begin{calculation}[A coordinate ring computed from a parametrized parabola]
This is a standard computation illustrating how an equation becomes a quotient
ring.  Let
\[
  X=\V(y-x^2)\subseteq \A^2_{\C}.
\]
The coordinate ring is
\[
  \C[X]=\C[x,y]/(y-x^2).
\]
Define a \(\C\)-algebra homomorphism
\[
  \Phi\colon \C[x,y]\longrightarrow \C[t],
  \qquad
  \Phi(x)=t,\quad \Phi(y)=t^2.
\]
Then
\[
  \Phi(y-x^2)=t^2-t^2=0,
\]
so \((y-x^2)\subseteq \ker\Phi\).  Conversely, every polynomial
\(f(x,y)\in \C[x,y]\) can be reduced modulo \(y-x^2\) by replacing \(y\) with
\(x^2\).  Thus
\[
  f(x,y)\equiv f(x,x^2)\pmod{(y-x^2)}.
\]
If \(\Phi(f)=f(t,t^2)=0\) as a polynomial in \(t\), then
\[
  f(x,x^2)=0,
\]
so the reduction of \(f\) is zero and \(f\in (y-x^2)\).  Hence
\[
  \ker\Phi=(y-x^2).
\]
By the first isomorphism theorem,
\[
  \C[x,y]/(y-x^2)\cong \C[t].
\]
Therefore
\[
  \K(X)=\Frac(\C[X])\cong \Frac(\C[t])=\C(t).
\]
In this example the coordinate \(t\) is the parameter, while the coordinate
functions on \(X\) satisfy \(x=t\) and \(y=t^2\).
\end{calculation}

\begin{calculation}[A second affine example]
Let
\[
  U=\V(xy-1)\subseteq \A^2_{\C}.
\]
The equation \(xy=1\) allows \(y=x^{-1}\).  Algebraically,
\[
  \C[U]=\C[x,y]/(xy-1).
\]
The homomorphism
\[
  \Psi\colon \C[x,y]\longrightarrow \C[t,t^{-1}],
  \qquad
  \Psi(x)=t,\quad \Psi(y)=t^{-1}
\]
satisfies \(\Psi(xy-1)=0\), and the same elimination argument gives
\[
  \C[x,y]/(xy-1)\cong \C[t,t^{-1}].
\]
Its function field is
\[
  \K(U)=\Frac(\C[t,t^{-1}])=\C(t).
\]
This example shows that different affine curves may have the same function
field even though their coordinate rings are not the same.
\end{calculation}

\begin{remark}[Section summary]
The coordinate ring records polynomial functions on a variety.  Passing to its
fraction field allows quotients of polynomial functions and produces the
function field.
\end{remark}

\sourceunclear{2}{The page contains illustrative local-chart drawings and a
quotient-like picture, but the precise equivalence relation and labels in
those drawings are not fully readable.}

\section{Projective Space and Projective Varieties}

Projective space is introduced because many naturally occurring curves are
best understood after adding their points at infinity.  This is also the
setting in which compact Riemann surfaces become algebraic curves: the
embedding theorem places them inside complex projective space.

% Source: PDF p. 3
\begin{definition}
Let \(k\) be a field and \(N\geq 0\).  Projective \(N\)-space over \(k\) is
\[
  \Pj^N_k
  =
  \bigl(k^{N+1}\setminus\{(0,\ldots,0)\}\bigr)\big/{\sim},
\]
where
\[
  (x_0,\ldots,x_N)\sim (y_0,\ldots,y_N)
\]
if and only if there exists \(\lambda\in k^\times\) such that
\[
  (y_0,\ldots,y_N)=(\lambda x_0,\ldots,\lambda x_N).
\]
The equivalence class of a nonzero vector is written
\[
  [x_0:\cdots:x_N].
\]
\end{definition}

\begin{remark}
The brackets \([x_0:\cdots:x_N]\) are homogeneous coordinates.  Multiplying
all coordinates by the same nonzero scalar does not change the point:
\[
  [x_0:\cdots:x_N]=[\lambda x_0:\cdots:\lambda x_N],
  \qquad \lambda\in k^\times.
\]
This is why equations on projective space must be homogeneous: if
\(f(x_0,\ldots,x_N)=0\), then the truth of the equation should not depend on
which representative of \([x_0:\cdots:x_N]\) is chosen.
\end{remark}

\begin{calculation}[Rescaling and homogeneous equations]
Let
\[
  F(X,Y,Z)=XZ-Y^2.
\]
This polynomial is homogeneous of degree \(2\).  If
\([X:Y:Z]=[\lambda X:\lambda Y:\lambda Z]\), then
\[
  F(\lambda X,\lambda Y,\lambda Z)
  =
  (\lambda X)(\lambda Z)-(\lambda Y)^2
  =
  \lambda^2(XZ-Y^2)
  =
  \lambda^2 F(X,Y,Z).
\]
Since \(\lambda\neq 0\), the equation \(F=0\) is true for one representative
if and only if it is true for every representative.  Therefore
\[
  XZ-Y^2=0
\]
is a well-defined equation on \(\Pj^2_{\C}\).

By contrast, the nonhomogeneous equation \(X+1=0\) is not well-defined on
\(\Pj^1_{\C}\).  The two vectors \((-1,1)\) and \((-2,2)\) represent the same
projective point:
\[
  [-1:1]=[-2:2].
\]
But
\[
  -1+1=0,\qquad -2+1=-1\neq 0.
\]
Thus the truth of the equation depends on the representative, so it cannot
define a subset of projective space.
\end{calculation}

\begin{definition}
A projective algebraic set in \(\Pj^N_k\) is the common zero locus of
homogeneous polynomials.  If \(I\subseteq k[X_0,\ldots,X_N]\) is a homogeneous
ideal, write
\[
  \V(I)=\{[x_0:\cdots:x_N]\in \Pj^N_k
  \mid f(x_0,\ldots,x_N)=0\text{ for every }f\in I\}.
\]
An irreducible projective algebraic set is a projective variety.
\end{definition}

\begin{remark}[Explanatory background]
Projective space is covered by affine charts.  For example, the chart
\(x_i\neq 0\) in \(\Pj^N_k\) is obtained by scaling so that \(x_i=1\).  This
is the projective analogue of choosing coordinates on a local patch.
\end{remark}

% Source: PDF p. 3
\begin{example}
The source lists several base fields that may be used in projective geometry:
\[
  \C,\qquad \Q,\qquad \F_q,\qquad \text{number fields}.
\]
\end{example}

\begin{remark}[Section summary]
Projective geometry uses homogeneous coordinates and homogeneous equations.
It is the natural setting for compact algebraic curves because it includes
points at infinity.
\end{remark}

\section{Realizing Smooth Manifolds as Varieties}

The source next asks when familiar smooth spaces can be described by
polynomial equations.  This is false for arbitrary smooth manifolds, but true
in the one-dimensional complex case that matters here.

% Source: PDF p. 3
\begin{remark}
Not every smooth manifold can be realized as an algebraic variety.  The source
records compact Riemann surfaces as an important positive case.
\end{remark}

% Source: PDF p. 3
\begin{theorem}[Compact Riemann surfaces as projective curves]
Let \(M\) be a compact Riemann surface.  Then \(M\) can be embedded into
complex projective space:
\[
  \varphi\colon M\hookrightarrow \Pj^N_{\C}
\]
for some \(N\).  Its image is a closed complex submanifold of
\(\Pj^N_{\C}\), and by Chow's theorem it is a closed algebraic subvariety.
Thus \(M\) can be regarded as a smooth projective algebraic curve over
\(\C\).
\end{theorem}

% Source: PDF p. 3
\begin{remark}
The conclusion is that, after choosing such an embedding,
\[
  M\cong \image(\varphi)\subseteq \Pj^N_{\C}.
\]
This is the bridge from the analytic object \(M\) to an algebraic object
defined by homogeneous equations.
\end{remark}

\begin{remark}[Section summary]
Compact Riemann surfaces are special: although arbitrary smooth manifolds need
not be algebraic, compact Riemann surfaces can be embedded as projective
algebraic curves over \(\C\).
\end{remark}

\sourceunclear{3}{The handwritten notes include an embedding diagram
\(M\to \Pj^N_{\C}\) and comments about injectivity/embedding.  The main meaning
is clear, but some side annotations are not readable enough to transcribe
exactly.}

\section{The Projective Line}

The simplest projective curve is \(\Pj^1_{\C}\).  It is the main comparison
object for the later elliptic curve: maps to \(\Pj^1_{\C}\) correspond to
subfields isomorphic to \(\C(x)\) inside a curve's function field.

% Source: PDF pp. 4--5
\begin{example}
The complex projective line is
\[
  \Pj^1_{\C}
  =
  \bigl(\C^2\setminus\{(0,0)\}\bigr)\big/{\C^\times}.
\]
It is identified with the Riemann sphere:
\[
  \Pj^1_{\C}\cong S^2.
\]
\end{example}

\begin{remark}[Affine charts on \(\Pj^1\)]
The two standard affine charts are obtained by choosing a nonzero coordinate.
If \(y\neq 0\), then
\[
  [x:y]=[x/y:1].
\]
If \(x\neq 0\), then
\[
  [x:y]=[1:y/x].
\]
Thus \(\Pj^1_{\C}\) is obtained by gluing two copies of the affine line.
\end{remark}

\begin{example}
On the chart \(y\neq 0\), the coordinate is \(t=x/y\).  On the chart
\(x\neq 0\), the coordinate is \(s=y/x\).  On the overlap both coordinates
are nonzero and satisfy
\[
  s=1/t.
\]
\end{example}

\begin{calculation}[The function field of \(\Pj^1_{\C}\)]
This is a standard computation behind the later notation
\(\K(\Pj^1_{\C})=\C(x)\).  On the affine chart \(y\neq 0\), write
\[
  t=x/y.
\]
Rational functions on this chart are elements of \(\C(t)\).  On the other
chart \(x\neq 0\), write
\[
  s=y/x.
\]
On the overlap, \(s=1/t\), so
\[
  \C(s)=\C(1/t)=\C(t).
\]
Thus the two affine charts give the same rational function field:
\[
  \K(\Pj^1_{\C})=\C(t).
\]
Later the same field is written \(\C(x)\) when the coordinate comes from the
\(x\)-coordinate map on the elliptic curve.
\end{calculation}

\begin{remark}[Section summary]
The curve \(\Pj^1_{\C}\) is the projective line, or Riemann sphere.  Its
function field is \(\C(x)\), so it becomes the base curve for the later
degree-two map from an elliptic curve.
\end{remark}

\section{Veronese Embeddings and Rational Curves}

The Veronese maps are the first concrete examples of parametrized projective
curves.  They show how a curve that is abstractly \(\Pj^1_{\C}\) can appear
inside a higher-dimensional projective space as the zero set of homogeneous
polynomial equations.

% Source: PDF p. 4
\begin{example}[The quadratic Veronese embedding]
Define
\[
  \nu_2\colon \Pj^1_{\C}\longrightarrow \Pj^2_{\C},
  \qquad
  [x:y]\longmapsto [x^2:xy:y^2].
\]
This is well-defined because replacing \((x,y)\) by \((\lambda x,\lambda y)\)
multiplies all three target coordinates by the common scalar \(\lambda^2\).
Writing homogeneous coordinates on \(\Pj^2_{\C}\) as \([X:Y:Z]\), every point
in the image satisfies
\[
  XZ-Y^2=0.
\]
Indeed,
\[
  (x^2)(y^2)-(xy)^2=0.
\]
\end{example}

\begin{proposition}
The map \(\nu_2\) is injective, and its image is the conic
\[
  \V(XZ-Y^2)\subseteq \Pj^2_{\C}.
\]
\end{proposition}

\begin{proof}
First suppose
\[
  \nu_2([x:y])=\nu_2([x':y']).
\]
If \(x\neq 0\), then the first and second coordinates of
\([x^2:xy:y^2]\) determine \(y/x\).  Hence
\([x:y]=[1:y/x]\) is determined by the image.  If \(x=0\), then the image is
\([0:0:1]\), so the preimage is \([0:1]\).  Thus \(\nu_2\) is injective.

We have already checked that \(\image(\nu_2)\subseteq \V(XZ-Y^2)\).  Conversely,
let \([X:Y:Z]\in \V(XZ-Y^2)\).  If \(X\neq 0\), then
\[
  [X:Y:Z]=[1:Y/X:Z/X],
\]
and the equation gives \(Z/X=(Y/X)^2\).  Hence
\[
  [X:Y:Z]=[1:t:t^2]=\nu_2([1:t])
  \qquad \text{with } t=Y/X.
\]
If \(X=0\), then the equation gives \(Y=0\), so the point is
\([0:0:1]=\nu_2([0:1])\).  Therefore the conic is exactly the image.
\end{proof}

\begin{calculation}[Recovering the parameter for the conic]
On the affine chart \(x\neq 0\) of the source, every point of
\(\Pj^1_{\C}\) has the form \([1:t]\).  Substituting into \(\nu_2\) gives
\[
  \nu_2([1:t])=[1:t:t^2].
\]
Thus on the target chart \(X\neq 0\), the parameter is recovered by
\[
  t=\frac{Y}{X},
  \qquad
  \frac{Z}{X}=t^2.
\]
The single source point with \(x=0\) is \([0:1]\), and its image is
\[
  \nu_2([0:1])=[0:0:1].
\]
This is the point at infinity on the conic relative to the affine chart
\(X\neq 0\).  The computation shows explicitly how the parametrization and
the equation \(XZ-Y^2=0\) encode the same curve.
\end{calculation}

% Source: PDF pp. 4--5
\begin{example}[The cubic Veronese embedding]
Define
\[
  \nu_3\colon \Pj^1_{\C}\longrightarrow \Pj^3_{\C},
  \qquad
  [x:y]\longmapsto [x^3:x^2y:xy^2:y^3].
\]
Writing homogeneous coordinates on \(\Pj^3_{\C}\) as \([X:Y:Z:W]\), the
source records the image as
\[
  \image(\nu_3)
  =
  \left\{
  [X:Y:Z:W]\in \Pj^3_{\C}
  \ \middle|\
  \begin{array}{l}
    Y^2=XZ,\\
    Z^2=YW,\\
    YZ=XW
  \end{array}
  \right\}.
\]
The equations follow by substituting
\[
  X=x^3,\qquad Y=x^2y,\qquad Z=xy^2,\qquad W=y^3.
\]
For example, \(Y^2=x^4y^2=XZ\), \(Z^2=x^2y^4=YW\), and
\(YZ=x^3y^3=XW\).
\end{example}

\begin{proposition}
The three equations displayed above cut out the image of \(\nu_3\).
\end{proposition}

\begin{proof}
The substitution in the previous example proves that every point of
\(\image(\nu_3)\) satisfies the three equations.  Conversely, let
\([X:Y:Z:W]\) satisfy them.

If \(X\neq 0\), scale the point so that \(X=1\).  Put \(t=Y/X\).  The equation
\(Y^2=XZ\) gives \(Z/X=t^2\), and the equation \(YZ=XW\) gives
\(W/X=t^3\).  Hence the point is
\[
  [1:t:t^2:t^3]=\nu_3([1:t]).
\]
If \(X=0\), then \(Y^2=XZ=0\), so \(Y=0\).  Then \(Z^2=YW=0\), so \(Z=0\).
The point is therefore \([0:0:0:1]=\nu_3([0:1])\).  Thus the common zero locus
is exactly \(\image(\nu_3)\).
\end{proof}

\begin{calculation}[Recovering the parameter for the twisted cubic]
On the source chart \(x\neq 0\), write a point as \([1:t]\).  Then
\[
  \nu_3([1:t])=[1:t:t^2:t^3].
\]
On the target chart \(X\neq 0\), the parameter is recovered from
\[
  t=\frac{Y}{X}.
\]
The defining equations then force the remaining coordinates:
\[
  \frac{Z}{X}
  =
  \left(\frac{Y}{X}\right)^2=t^2,
  \qquad
  \frac{W}{X}
  =
  \frac{YZ}{X^2}=t^3.
\]
At the source point \([0:1]\), the image is
\[
  \nu_3([0:1])=[0:0:0:1].
\]
This is the point at infinity of the twisted cubic with respect to the chart
\(X\neq 0\).
\end{calculation}

\begin{remark}
The curves \(\image(\nu_2)\) and \(\image(\nu_3)\) are rational curves because
they are parametrized by \(\Pj^1_{\C}\).  The quadratic case gives a conic in
\(\Pj^2_{\C}\), while the cubic case gives a twisted cubic in
\(\Pj^3_{\C}\).
\end{remark}

\begin{remark}[Section summary]
The Veronese maps show how a parametrized curve can be translated into
implicit homogeneous equations.  This prepares the way for studying curves by
their equations rather than by pictures alone.
\end{remark}

\section{A Plane Cubic Model of an Elliptic Curve}

The rational curves above have function field \(\C(t)\).  The next example is
different: a nonsingular plane cubic is a genus-one curve.  Analytically it is
related to a complex torus, while algebraically it is described by a
Weierstrass equation.

% Source: PDF pp. 6--7
\begin{example}
The source records a map from a torus \(T\) into \(\Pj^2_{\C}\) whose image is
given by the homogeneous cubic equation
\[
  Y^2Z-X^3-aXZ^2-bZ^3=0.
\]
Equivalently,
\[
  Y^2Z=X^3+aXZ^2+bZ^3.
\]
The corresponding projective curve is
\[
  E=\V(Y^2Z-X^3-aXZ^2-bZ^3)\subseteq \Pj^2_{\C}.
\]
When this cubic is nonsingular, it is an elliptic curve.
\end{example}

\begin{remark}
For the Weierstrass cubic above, nonsingularity is equivalent to the cubic
polynomial
\[
  x^3+ax+b
\]
having distinct roots.  Equivalently, its discriminant is nonzero:
\[
  -16(4a^3+27b^2)\neq 0.
\]
This condition is not written explicitly in the source, but it is the standard
condition needed for the phrase ``elliptic curve'' in this model.
\end{remark}

% Source: PDF p. 7
\begin{remark}[The affine chart \(Z\neq 0\)]
On the inhomogeneous open set \(Z\neq 0\), choose the representative with
\(Z=1\).  Put
\[
  x=X/Z,\qquad y=Y/Z.
\]
Then the projective equation becomes the affine Weierstrass equation
\[
  y^2=x^3+ax+b.
\]
Thus the affine part of \(E\) is
\[
  E\cap\{Z\neq 0\}
  =
  \{(x,y)\in \C^2\mid y^2=x^3+ax+b\}.
\]
\end{remark}

\begin{calculation}[Dehomogenizing the cubic]
On the chart \(Z\neq 0\), divide
\[
  Y^2Z=X^3+aXZ^2+bZ^3
\]
by \(Z^3\).  This gives
\[
  \left(\frac{Y}{Z}\right)^2
  =
  \left(\frac{X}{Z}\right)^3
  +a\left(\frac{X}{Z}\right)+b.
\]
With
\[
  x=X/Z,\qquad y=Y/Z,
\]
the equation becomes
\[
  y^2=x^3+ax+b.
\]
Conversely, any affine solution \((x,y)\) gives the projective point
\[
  [X:Y:Z]=[x:y:1],
\]
and direct substitution gives
\[
  y^2\cdot 1=x^3+a x\cdot 1^2+b\cdot 1^3.
\]
Thus the affine equation and the \(Z\neq 0\) part of the projective equation
are exactly the same calculation in different coordinates.
\end{calculation}

\begin{remark}[The point at infinity]
The points of \(E\) with \(Z=0\) satisfy
\[
  -X^3=0,
\]
so \(X=0\).  Since \([0:0:0]\) is not a point of projective space, \(Y\neq 0\),
and there is exactly one point at infinity:
\[
  [0:1:0].
\]
Therefore
\[
  E=
  \{(x,y)\in \C^2\mid y^2=x^3+ax+b\}\cup\{[0:1:0]\}.
\]
\end{remark}

\begin{calculation}[The point at infinity and the sign involution]
Set \(Z=0\) in the homogeneous equation:
\[
  Y^2Z-X^3-aXZ^2-bZ^3=0.
\]
Every term containing \(Z\) vanishes, leaving
\[
  -X^3=0.
\]
Hence \(X=0\).  A projective point cannot have all coordinates zero, so
\(Y\neq 0\).  Rescaling by \(Y^{-1}\) gives the unique point
\[
  [0:Y:0]=[0:1:0].
\]

On the affine chart, the map
\[
  \iota(x,y)=(x,-y)
\]
preserves the equation because
\[
  (-y)^2=y^2=x^3+ax+b.
\]
The point at infinity is also fixed by the corresponding projective map
\[
  [X:Y:Z]\longmapsto [X:-Y:Z],
\]
since
\[
  [0:-1:0]=[0:1:0].
\]
\end{calculation}

\begin{remark}[Section summary]
The projective cubic contains the affine Weierstrass curve together with one
point at infinity.  When nonsingular, this projective curve is an elliptic
curve.
\end{remark}

\begin{remark}[The 180-degree quotient picture]
% Source: PDF p. 6
The source also draws a torus with a \(180^\circ\) rotation and a folding
picture.  This should be read as a quotient picture, not as a deformation that
turns a torus itself into a sphere.  Topologically, a complex elliptic curve is
a torus.  The order-two symmetry may be visualized as a \(180^\circ\) rotation
identifying paired points.  In the Weierstrass model, the same symmetry is the
involution
\[
  \iota\colon E\longrightarrow E,
  \qquad
  \iota(x,y)=(x,-y).
\]
Taking the quotient means that a point and its partner under \(\iota\) are
regarded as the same point.  The quotient by this involution is the projective
line \(\Pj^1_{\C}\), which is the Riemann sphere.  Algebraically, the quotient
map is the two-to-one map given by the \(x\)-coordinate, developed in the next
section.
\end{remark}

\sourceunclear{6}{The source draws a map from a torus \(T\) to the plane
cubic, but it does not specify the formula for the map.  These notes therefore
record only the image equation.  The \(180^\circ\) quotient picture is
included above, but the source does not provide its full branch or fixed-point
data.}

\section{The Function Field of the Elliptic Curve}

The field-theoretic content of the elliptic-curve example appears when we
compare the function field of \(E\) with the function field of \(\Pj^1_{\C}\).
The \(x\)-coordinate gives a two-to-one map \(E\to\Pj^1_{\C}\), and this
becomes a degree-two extension of function fields.

% Source: PDF pp. 8--9
\begin{definition}
Let
\[
  E=\V(Y^2Z-X^3-aXZ^2-bZ^3)\subseteq \Pj^2_{\C}
\]
be the nonsingular projective cubic above.  As standard background, passing
from an irreducible curve to a nonempty affine open subset does not change its
function field.  On the affine chart \(Z\neq 0\), the coordinate ring is
\[
  \C[x,y]\big/(y^2-x^3-ax-b).
\]
Its function field is the fraction field
\[
  \K(E)
  =
  \Frac\!\left(\C[x,y]\big/(y^2-x^3-ax-b)\right).
\]
Equivalently,
\[
  \K(E)\cong \C(x)[T]\big/(T^2-x^3-ax-b),
\]
where the class of \(T\) is the function \(y\).
\end{definition}

\begin{calculation}[Deriving \(\K(E)=\C(x)(y)\)]
Let
\[
  f(x)=x^3+ax+b.
\]
On the affine chart \(Z\neq 0\), the coordinate ring is
\[
  A=\C[x,y]/(y^2-f(x)).
\]
No nonzero polynomial in \(\C[x]\) lies in the ideal \((y^2-f(x))\), so the
class of \(x\) gives a copy of \(\C[x]\) inside \(A\).  In particular, inside
the fraction field of \(A\), the element \(x\) is transcendental over \(\C\)
and \(y\) satisfies
\[
  y^2=f(x).
\]
After passing from \(\C[x]\) to its fraction field \(\C(x)\), we obtain the
localized algebra
\[
  \C(x)[y]/(y^2-f(x)).
\]
Every polynomial in \(y\) can be reduced to degree at most one using
\(y^2=f(x)\).  Explicitly,
\[
  y^{2m}=f(x)^m,
  \qquad
  y^{2m+1}=f(x)^m y.
\]
Therefore every element of \(\C(x)[y]/(y^2-f(x))\) has the form
\[
  a(x)+b(x)y,
  \qquad a(x),b(x)\in \C(x).
\]
Since \(f(x)\) has distinct roots, it is not a square in \(\C(x)\).  Hence
\(T^2-f(x)\) is irreducible over \(\C(x)\), so this localized algebra is a
field.  It is therefore the fraction field of \(A\), and
\[
  \K(E)=\C(x)(y),
  \qquad y^2=f(x).
\]
Thus \(\K(E)\) is generated over \(\C(x)\) by one square root, and every
element of \(\K(E)\) can be written uniquely as \(a(x)+b(x)y\).
\end{calculation}

% Source: PDF pp. 8--9
\begin{remark}
The function field of the projective line is
\[
  \K(\Pj^1_{\C})=\C(x).
\]
The rational function \(x=X/Z\) on \(E\) has a pole at the point at infinity.
By the standard fact that a nonconstant rational function on a nonsingular
projective curve defines a morphism to \(\Pj^1\), it gives a morphism
\[
  E\longrightarrow \Pj^1_{\C}.
\]
On affine points this morphism is simply
\[
  (x,y)\longmapsto [x:1],
\]
and the point at infinity \([0:1:0]\) maps to \([1:0]\).
At the level of function fields, it corresponds to the inclusion
\[
  \C(x)\subseteq \K(E).
\]
This is an example of the general reversal: a map of varieties gives an
inclusion of function fields in the opposite direction.
\end{remark}

\[
\begin{tikzcd}
  E \arrow[r, "\pi=x"] & E/\langle\iota\rangle \cong \Pj^1_{\C}\\
  \K(E) & \K(E)^{\langle\iota^*\rangle}=\K(\Pj^1_{\C})=\C(x) \arrow[l, hook]
\end{tikzcd}
\]

% Source: PDF p. 8
\begin{remark}[The degree-two map]
For a generic value of \(x\), the equation
\[
  y^2=x^3+ax+b
\]
has two solutions:
\[
  y=\pm\sqrt{x^3+ax+b}.
\]
Therefore the morphism \(E\to \Pj^1_{\C}\) defined by \(x\) is a degree-two
map.  The source records this as a ``\(2:1\) map.''
\end{remark}

% Source: PDF p. 9
\begin{proposition}
Let
\[
  F=\C(x),
  \qquad
  L=\K(E).
\]
Then
\[
  L=F(y),\qquad y^2=x^3+ax+b,
\]
and \(L/F\) is a quadratic extension.  Moreover,
\[
  \Gal(L/F)\cong \Z/2\Z.
\]
\end{proposition}

\begin{proof}
Since \(L\cong \C(x)[T]/(T^2-x^3-ax-b)\), the field \(L\) is obtained from
\(F=\C(x)\) by adjoining a root \(y\) of
\[
  T^2-(x^3+ax+b)\in F[T].
\]
Because \(x^3+ax+b\) has distinct roots, it has simple zeros as a rational
function of \(x\).  A square in \(\C(x)\) has every zero and pole of even
order, so \(x^3+ax+b\) is not a square in \(F\).  Therefore
\[
  [L:F]=2.
\]
The nontrivial automorphism over \(F\) fixes \(x\) and sends
\[
  y\longmapsto -y.
\]
Indeed, this assignment preserves the defining relation because
\[
  (-y)^2=y^2=x^3+ax+b.
\]
It has order two, and a quadratic extension has at most two \(F\)-automorphisms.
Thus
\[
  \Gal(L/F)=\{1,\sigma\},
  \qquad
  \sigma(x)=x,\quad \sigma(y)=-y,
\]
so
\[
  \Gal(L/F)\cong \Z/2\Z.
\]
\end{proof}

\begin{remark}
Geometrically, the involution \(\iota(x,y)=(x,-y)\) exchanges the two points in
a generic fiber of the map \(E\to \Pj^1_{\C}\).  Algebraically, the induced
automorphism \(\iota^*\) is the nontrivial automorphism of the quadratic
function-field extension \(\K(E)/\C(x)\).
\end{remark}

\begin{calculation}[Computing the invariant field]
Write a general element of \(\K(E)\) over \(\C(x)\) as
\[
  u=a(x)+b(x)y,
  \qquad a(x),b(x)\in \C(x).
\]
The pullback action of the involution is
\[
  \iota^*(x)=x,\qquad \iota^*(y)=-y.
\]
Therefore
\[
  \iota^*(u)
  =
  \iota^*\bigl(a(x)+b(x)y\bigr)
  =
  a(x)-b(x)y.
\]
The element \(u\) is invariant precisely when
\[
  a(x)+b(x)y=a(x)-b(x)y.
\]
Subtracting gives
\[
  2b(x)y=0.
\]
Since the ground field is \(\C\), the characteristic is zero, and \(y\neq 0\)
in the function field.  Hence \(b(x)=0\).  Therefore the invariant elements
are exactly
\[
  u=a(x)\in \C(x),
\]
and so
\[
  \K(E)^{\langle \iota^*\rangle}=\C(x).
\]
\end{calculation}

\begin{remark}[Invariant functions]
In standard algebraic language, checking which quantities are unchanged by the
\(180^\circ\) symmetry means taking invariants.  The induced action on rational
functions is
\[
  \iota^*\colon \K(E)\longrightarrow \K(E),
  \qquad
  \iota^*(x)=x,\quad \iota^*(y)=-y.
\]
Then a rational function \(f\in \K(E)\) is unchanged by the symmetry precisely
when
\[
  \iota^*(f)=f.
\]
The fixed field of this order-two action is
\[
  \K(E)^{\langle\iota^*\rangle}=\C(x).
\]
An invariant rational function has the same value at paired points
\[
  P\quad\text{and}\quad \iota(P)
\]
wherever both values are defined, so it descends to a rational function on the
quotient curve.  Thus the invariant functions are exactly the rational
functions of the \(x\)-coordinate, which is why the quotient curve has function
field \(\C(x)=\K(\Pj^1_{\C})\).
\end{remark}

\begin{remark}[Section summary]
The elliptic curve function field is obtained from \(\C(x)\) by adjoining
one square root.  The sign change of that square root is the nontrivial Galois
automorphism, and the functions fixed by this sign change are exactly
\(\C(x)\).
\end{remark}

\sourceunclear{8}{The page indicates a two-to-one map and sketches the square
root description of the fibers.  Branch points or ramification data are not
specified in the source.}

\section{Summary: From Geometry to Galois Theory}

The lecture's main chain of ideas can now be read in one line:
\[
  E:\ y^2=x^3+ax+b
  \quad\leadsto\quad
  \K(E)=\C(x)(y)
  \quad\leadsto\quad
  \K(E)/\C(x).
\]
The map \(E\to\Pj^1_{\C}\) defined by the \(x\)-coordinate is generically
two-to-one, because a typical value of \(x\) gives the two points
\[
  (x,\sqrt{x^3+ax+b})
  \qquad\text{and}\qquad
  (x,-\sqrt{x^3+ax+b}).
\]
On function fields, this means that \(\K(E)\) is obtained from \(\C(x)\) by
adjoining one square root:
\[
  y^2=x^3+ax+b.
\]
The nontrivial symmetry of the extension fixes \(x\) and sends \(y\) to
\(-y\).  Therefore the Galois group is the two-element group
\[
  \Gal(\K(E)/\C(x))\cong \Z/2\Z.
\]
Equivalently, the base field \(\C(x)\) is the field of rational functions on
the quotient \(E/\langle\iota\rangle\cong \Pj^1_{\C}\).

\section{Exercises}

\begin{exercise}
Let \(F\subseteq E\) be a finite field extension.  Explain explicitly why
\(E\) is an \(F\)-vector space, and identify the scalar multiplication map.
\end{exercise}

\begin{exercise}
For the extension \(\Q(\sqrt d)/\Q\), with \(d\) squarefree and not a square,
write the matrix for multiplication by \(a+b\sqrt d\) in the basis
\((1,\sqrt d)\).
\end{exercise}

\begin{exercise}
Compute the coordinate ring and function field of the affine parabola
\[
  X=\V(y-x^2)\subseteq \A^2_{\C}.
\]
Write the homomorphism to \(\C[t]\), compute its kernel, and identify the
fraction field.
\end{exercise}

\begin{exercise}
Compute the coordinate ring and function field of
\[
  U=\V(xy-1)\subseteq \A^2_{\C}.
\]
Explain why \(y\) becomes \(x^{-1}\) in the coordinate ring.
\end{exercise}

\begin{exercise}
Show that the equation \(XZ-Y^2=0\) is well-defined on \(\Pj^2_{\C}\); that
is, its truth is unchanged when \([X:Y:Z]\) is rescaled.
\end{exercise}

\begin{exercise}
Give an explicit example showing that the equation \(X+1=0\) is not
well-defined on \(\Pj^1_{\C}\).
\end{exercise}

\begin{exercise}
For the quadratic Veronese map
\[
  \nu_2([x:y])=[x^2:xy:y^2],
\]
verify by substitution that every image point satisfies \(XZ-Y^2=0\).  Then
recover the parameter \(t\) on the chart \(X\neq 0\).
\end{exercise}

\begin{exercise}
Repeat the injectivity argument for the cubic Veronese map
\[
  \nu_3([x:y])=[x^3:x^2y:xy^2:y^3].
\]
\end{exercise}

\begin{exercise}
Verify directly that the three equations
\[
  Y^2=XZ,\qquad Z^2=YW,\qquad YZ=XW
\]
hold for every point in the image of \(\nu_3\).
\end{exercise}

\begin{exercise}
On the chart \(Z\neq 0\), check that the homogeneous equation
\[
  Y^2Z=X^3+aXZ^2+bZ^3
\]
reduces to \(y^2=x^3+ax+b\).  Then check separately that the only point of
the curve with \(Z=0\) is \([0:1:0]\).
\end{exercise}

\begin{exercise}
Verify directly that the transformation
\[
  (x,y)\longmapsto (x,-y)
\]
preserves the affine equation \(y^2=x^3+ax+b\).
\end{exercise}

\begin{exercise}
Let \(f(x)=x^3+ax+b\).  In the field \(\C(x)(y)\) with \(y^2=f(x)\), reduce
\[
  c_0(x)+c_1(x)y+c_2(x)y^2+c_3(x)y^3
\]
to the form \(A(x)+B(x)y\), where \(A(x),B(x)\in \C(x)\).
\end{exercise}

\begin{exercise}
Let \(\iota^*\) be the automorphism of \(\K(E)\) determined by
\(\iota^*(x)=x\) and \(\iota^*(y)=-y\).  Using
\[
  \K(E)=\C(x)\oplus \C(x)y
\]
as a vector space over \(\C(x)\), show that the invariant subfield is
\[
  \K(E)^{\langle\iota^*\rangle}=\C(x).
\]
\end{exercise}

\begin{exercise}
Let \(L=\C(x)(y)\), where \(y^2=x^3+ax+b\), and assume
\(x^3+ax+b\) has distinct roots.  Verify that the map
\[
  \sigma(x)=x,\qquad \sigma(y)=-y
\]
defines a nontrivial \(\C(x)\)-automorphism of \(L\), and compute
\(\sigma^2\).
\end{exercise}

\begin{exercise}
Explain in words why invariant rational functions are constant on paired
points \(P\) and \(\iota(P)\), and why this means they descend to rational
functions on the quotient \(E/\langle\iota\rangle\).
\end{exercise}

\section*{Hints}

\begin{enumerate}
\item Use the field operations of \(E\): addition is addition in \(E\), and
scalar multiplication is \((a,v)\mapsto av\) for \(a\in F\), \(v\in E\).
\item Compute
\((a+b\sqrt d)(r+s\sqrt d)=(ar+bds)+(as+br)\sqrt d\).
\item Use the map \(x\mapsto t\), \(y\mapsto t^2\).  The relation
\(y=x^2\) lets you eliminate \(y\), so the coordinate ring is isomorphic to
\(\C[t]\).
\item Use the map \(x\mapsto t\), \(y\mapsto t^{-1}\).  The relation
\(xy=1\) makes \(y\) the inverse of \(x\).
\item Substitute \((\lambda X,\lambda Y,\lambda Z)\) into \(XZ-Y^2\).
\item Compare the two representatives \([-1:1]\) and \([-2:2]\).
\item Substitute \(X=x^2\), \(Y=xy\), and \(Z=y^2\).  On \(X\neq 0\), use
\(t=Y/X\).
\item Work on the charts \(x\neq 0\) and \(x=0\), as in the proof for
\(\nu_2\).
\item Substitute \(X=x^3\), \(Y=x^2y\), \(Z=xy^2\), and \(W=y^3\).
\item Set \(Z=1\) on the affine chart.  For \(Z=0\), the homogeneous equation
forces \(X=0\), leaving the single projective point \([0:1:0]\).
\item The square of \(-y\) is the same as the square of \(y\).
\item Replace \(y^2\) by \(f(x)\) and \(y^3\) by \(f(x)y\).
\item Write a general element as \(a(x)+b(x)y\).  Applying \(\iota^*\) gives
\(a(x)-b(x)y\), so invariance forces \(b(x)=0\).
\item Check that \(\sigma\) preserves the relation \(y^2=x^3+ax+b\), then
apply \(\sigma\) twice.
\item If \(f(P)=f(\iota(P))\), then the value depends only on the equivalence
class of \(P\) in the quotient.
\end{enumerate}

\section*{Further Reading}

The references below are standard background sources and are not part of the
handwritten lecture source.
\begin{itemize}
\item Robin Hartshorne, \emph{Algebraic Geometry}, Chapters I.1--I.3, for
varieties, projective space, and morphisms.
\item Igor Shafarevich, \emph{Basic Algebraic Geometry 1}, for an introductory
treatment of affine and projective varieties.
\item William Fulton, \emph{Algebraic Curves}, for plane curves and function
fields.
\item Rick Miranda, \emph{Algebraic Curves and Riemann Surfaces}, for the
bridge between compact Riemann surfaces and algebraic curves.
\item Joseph H. Silverman and John Tate, \emph{Rational Points on Elliptic
Curves}, for elliptic curves in Weierstrass form.
\end{itemize}

\section*{Unclear Source Items}

\begin{itemize}
\item TODO: unclear from source, page 1.  Several arrows in the operation
diagram for \((F,+_F,\cdot_F)\) and \((E,+_E,\cdot_E)\) are illegible.
\item TODO: unclear from source, page 2.  The illustrative quotient/local-chart
drawings are not labeled precisely enough to reconstruct as formal diagrams.
\item TODO: unclear from source, page 3.  Some side annotations near the
embedding \(M\to\Pj^N_{\C}\) are not readable.
\item TODO: unclear from source, page 6.  The explicit torus-to-cubic map is
not given.
\item TODO: unclear from source, page 8.  The square-root sketch does not
include branch or ramification details.
\end{itemize}

\end{document}
